\section{第~\ref{sec:chemoutput}~章　化学输出}

通往心脏的两根导线及其不同的速度、激素级联中的负反馈、剂量频率编码、昼夜节律发生器及其光校准，都有直接测量；界限 $K<8$ 是本章推导的控制理论定理；由级联反馈本身产生的脉动是一个有强有力实验支持的假说。

{\footnotesize
\setlength{\tabcolsep}{3pt}\renewcommand{\arraystretch}{1.15}
\begin{longtable}{>{\raggedright}p{0.5cm}>{\raggedright}p{4.0cm}>{\raggedright}p{5.6cm}>{\raggedright}p{2.3cm}>{\raggedright\arraybackslash}p{1.7cm}}
\toprule
编号 & 论断 & 实验及测量内容 & 来源 & 状态 \\
\midrule\endfirsthead
\toprule
编号 & 论断 & 实验及测量内容 & 来源 & 状态 \\
\midrule\endhead
1 & 心脏起搏器自身每分钟跳动约 $100$ 次；静息时刹车踩着，频率通常约为 $60$--$70$ & 人体，完全阻断通往心脏的两根导线：固有心率高于静息心率，并随年龄下降，成人约为每分钟 $100$ 次。可见静息时刹车占主导。静息心率 $60$--$70$ 是典型值，未单独引用。 & \cite{JoseCollison1970ihr} & 已测量 \\
2 & 油门 \nt{NORA} 和刹车 \nt{ACET} 是低通滤波器，刹车更快~\eqref{eq:gasbrake} & 狗，对心脏迷走神经或交感神经进行调频刺激，测量到心房频率的传递函数：两条通路都是低通滤波器，交感通路的截止频率低得多，并附加约 $1.7\,\text{s}$ 的纯延迟。特性依赖于平均张力，因此线性公式~\eqref{eq:gasbrake}~只是近似。 & \cite{Berger1989} & 已测量 \\
3 & 刹车之所以快，是因为 \nt{ACET} 不经第二信使直接打开钾通道，而 \nt{NORA} 要经过一连串反应 & 心房细胞膜片：乙酰胆碱开启的钾通道由与受体偶联的 G 蛋白 $\beta\gamma$ 亚基打开，不需要细胞内的第二信使。\nt{NORA} 经 cAMP 的通路此处未引用。 & \cite{Logothetis1987gbg} & 已测量 \\
4 & 血液绕身体一周约需一分钟；血流中的 \nt{ADRE} 到达所有带受体的器官 & 数量级的一般估计；本章即如此表述，未引用来源。 & --- & 估计 \\
5 & 激素级联被负反馈包围：终端激素抑制前面几级 & 实验综述：肾上腺皮质激素在快、中、慢三个时间范围内抑制垂体释放 ACTH 和下丘脑释放释放激素。 & \cite{KellerWood1984fb} & 已测量 \\
6 & 三个相同的级在 $K<8$ 时稳定~\eqref{eq:cascadestable}，临界处周期为 $2\pi\tau/\sqrt3\approx3.6\tau$ & 本章推导：在频率 $\sqrt3/\tau$ 处，三级给出 $180^\circ$ 相移和 $8$ 倍衰减。 & 本章推导 & 理论 \\
7 & 电平误差为 $1/(1+K)$，因此这种调节器的精度不会优于约 $10\,\%$（命题~\ref{prop:cascade}） & 第~6 行对比例反馈的推论。真实的轴在多个级上、在多个时间范围内都有反馈（第~5 行），因此这一界限属于模型~\eqref{eq:cascade}，并非测量所得。 & 本章推导 & 理论 \\
8 & $K=4$ 和 $7$ 时电平趋于稳定，$K=12$ 时出现周期为 $3.7$--$3.9\,\tau$ 的平稳脉动（图~\ref{fig:cascade}） & 用 \texttt{models/\allowbreak chem\_\allowbreak models.py} 计算：$K=12$ 时相继的周期为 $3.9$、$3.7$、$3.8$、$3.7\,\tau$；$K=4$ 时计算结束时的幅度为 $0.003$。 & \texttt{models/\allowbreak chem\_\allowbreak models.py} & 模型 \\
9 & 应激激素大约每小时释放一次脉冲；脉冲由反馈本身产生，无需单独的发生器 & 大鼠，持续输注释放激素：ACTH 和皮质酮以接近每小时一次的生理频率振荡——脉冲并不需要来自下丘脑的节律性输入。带延迟反馈的垂体—肾上腺模型能给出这样的振荡。 & \cite{Walker2012glc, Walker2010} & 假说 \\
10 & 接收方读取的是剂量频率，而不是水平 & 自身不能释放促性腺激素释放激素的猴子：持续输注不能恢复垂体激素的释放，每小时一次的剂量能恢复；改回持续输注后响应再次消失。受体疲劳相当于高通滤波器是一种解释。 & \cite{Belchetz1978} & 已测量 \\
11 & 不同的剂量频率对同一腺体的不同细胞作用不同 & 同一批猴子：剂量频率提高到每小时 $2$--$5$ 次时，两种垂体激素都下降；降低到每 $3$~h 一次时，一种（LH）下降而另一种（FSH）上升，因此两者的比值由频率决定。 & \cite{Wildt1981gnrh} & 已测量 \\
12 & 昼夜节律由细胞内延迟数小时的负反馈产生 & 综述：时钟基因的蛋白质在延迟后抑制自身基因的转录；由转录和翻译构成的环路给出约一天的周期。 & \cite{TakahashiJS2017} & 已测量 \\
13 & 主发生器是几万个神经元组成的一群细胞，它使身体其余部分同步 & 综述：视交叉上核是主昼夜节律发生器；培养中的单个神经元就是独立的振荡器，网络使它们同步；小鼠每侧约有 $10^4$ 个神经元。 & \cite{Welsh2010scn} & 已测量 \\
14 & 人体的固有周期为 $24.18$~h & 在昏暗光线下、按脱离昼夜的作息生活的人：褪黑素、体温和皮质醇节律的周期在年轻人和老年人中平均均为 $24.18$~h，离散很小。 & \cite{Czeisler1999} & 已测量 \\
15 & 光像锁相环一样校准发生器~\eqref{eq:pll}；每天需要移动约 $11$~min & 人体，在一天中不同时间给予一次 $6.7$~h 的强光：得到幅度为 $5$~h 的第一类相位响应曲线——体温最低点之前给光造成延迟，之后给光造成提前。方程~\eqref{eq:pll}~是标准模型；$11$~min $=0.18$~h 是算术。 & \cite{Khalsa2003prc} & 已测量 \\
16 & 没有光就没有校准，节律漂移到固有周期 & 完全失明、没有光感、在正常社会中生活的人：约有一半人的褪黑素和皮质醇节律按自己的周期运行，与昼夜不一致。 & \cite{Sack1992blind} & 已测量 \\
\bottomrule
\end{longtable}}
